\clearpage
\section{Global RRP 的凸优化结构}
\subsection{符号与决策时点}
设 $t$ 为调仓日，$\mathcal A_t$ 为当期合格资产集合，$n_t$ 为其大小。历史数据严格满足日期早于 $t$。令 $\mu_t$ 为年化收益估计，$\Sigma_t$ 为年化协方差，$q_t$ 为风险预算参考，$w_t$ 为最终资本权重。所有优化变量只对 $\mathcal A_t$ 定义，再按原资产池顺序扩展回 30 维向量。不合格资产对应权重为零。

这些符号具有明确的时间含义。$q_t$ 和 $w_t$ 都是利用同一信息集合求得的目标，$w_t^-$ 则是实际交易前漂移权重。将 $w_t^-$ 与上一期目标直接混同，会忽略价格变化造成的自然权重漂移，并改变交易成本。收益预测 $\mu_t$ 也不等于调仓日实现收益，二者在程序中处于不同阶段。

\subsection{风险贡献分解}
对正波动组合，定义
\[
\sigma(w)=\sqrt{w^\mathsf T\Sigma_t w},\qquad
\frac{\partial\sigma(w)}{\partial w_i}=\frac{(\Sigma_t w)_i}{\sigma(w)}.
\]
资产的风险贡献为
\[
RC_i(w)=w_i\frac{(\Sigma_t w)_i}{\sigma(w)}.
\]
将各项相加，利用二次型定义可得
\[
\sum_i RC_i(w)=\frac{w^\mathsf T\Sigma_t w}{\sigma(w)}=\sigma(w).
\]
这一恒等式解释了风险贡献与总波动之间的联系。它并未要求各项必须非负，也未要求各项相等。某些资产与组合的协方差可能为负，因此负边际贡献在一般组合中具有数学可能性。风险预算问题通过选择特定参考权重，将贡献比例与正预算相匹配。

令预算 $b_i=1/n_t$，理想参考满足
\[
q_i(\Sigma_t q)_i=b_i q^\mathsf T\Sigma_t q.
\]
直接把这些乘积等式作为最终问题的约束，并不能由其形式立即保证凸性。本文改为求解可验证的对数风险预算问题，再将其归一化结果作为参考。

\subsection{对数风险预算参考}
在 $x_i>0$ 上求解
\[
\min_x f_t(x)=\frac12x^\mathsf T\Sigma_t x-\sum_i b_i\log x_i.
\]
其梯度与 Hessian 分别为
\[
\nabla f_t(x)=\Sigma_t x-(b_i/x_i)_i,
\]
\[
\nabla^2 f_t(x)=\Sigma_t+\mathop{\rm diag}(b_i/x_i^2).
\]
当协方差半正定且所有 $b_i>0$ 时，Hessian 在正正交域内正定，因而目标严格凸。若协方差正定，二次项保证目标在远离原点时增长，负对数项阻止分量趋于零，可以得到有限内部解。实际程序对协方差施加正特征值下限，相关条件需要与数值修复记录一起解释。

内部最优点满足 $x_i(\Sigma_t x)_i=b_i$。归一化 $q=x/(\mathbf1^\mathsf T x)$ 后，有
\[
q_i(\Sigma_t q)_i=\frac{b_i}{(\mathbf1^\mathsf T x)^2},\qquad
q^\mathsf T\Sigma_t q=\frac{1}{(\mathbf1^\mathsf T x)^2}.
\]
因此各资产的方差贡献比例等于预算。这一推导表明，归一化不是任意的后处理，而是与风险贡献的齐次结构相容。需要注意，等式对应求解时使用的协方差，不保证未来实现贡献仍然相等。

\subsection{最终权重问题}
指定 Global RRP 求解
\[
\min_{w\in\Delta_t}J_t(w)
=\|w-q_t\|_2^2
+\lambda_{v,t}\frac{w^\mathsf T\Sigma_t w}{v_t}
+\lambda_{r,t}\left[\max\left(0,\frac{R_t-\mu_t^\mathsf T w}{s_t}\right)\right]^2,
\]
\[
\Delta_t=\{w\in\mathbb R^{n_t}:\mathbf1^\mathsf T w=1,\ w\geq0\}.
\]
收益目标取当期风险预算参考的预测收益，
\[
R_t=\mu_t^\mathsf Tq_t.
\]
参考权重 $q_t$ 属于单纯形，因此该目标在当期估计下可由参考组合达到。它不再依赖固定收益倍数，也不等于对随后实现收益作保证。

归一化尺度为
\[
v_t=\max(e_t^\mathsf T\Sigma_t e_t,\varepsilon_v),\quad
e_t=\frac1{n_t}\mathbf1,
\]
\[
s_t=\max(|R_t|,\max_i|\mu_{i,t}|,\varepsilon_m).
\]
其中 $\varepsilon_v$ 为协方差数值修复使用的正尺度下限，$\varepsilon_m$ 为机器精度下限。二者只防止除零，不代表投资者的收益或风险要求。

\subsection{年度参数校准}
两项惩罚系数按年度更新。设某一年度的生效日为 $d_y$。程序只使用 $d_y$ 之前的数据，并依次截取两个连续的 252 日区间。较早区间用于确定候选尺度，随后区间用于候选验证。

在较早区间中，先以单位惩罚运行逐期优化，记录参考跟踪、归一化方差和归一化收益短缺三项的实际值。方差惩罚的三个候选取跟踪项与方差项比值的 25\%、50\% 和 75\% 分位数，收益短缺惩罚采用同样规则。两个三元素集合形成九组预先确定的组合，候选范围随历史目标项的量级变化。

随后 252 日区间按扣费净收益计算夏普。程序保留距离最高夏普不超过一倍标准误的候选，并优先选择在网格中存在相邻合格点的配置。若仍有多个候选，则先比较年化换手，再比较季度夏普中位数，最后按较小惩罚取值排序。选定参数从 $d_y$ 起冻结，直到下一年度更新。

这套程序限定了参数来源和信息时间，但没有使系数获得唯一统计识别。十个年度中，每年九组候选都进入一倍标准误集合，能够排除候选的年度数为零。年度选择主要由换手和稳定性规则决定，因此表中的系数应解释为可复现决策规则的输出，不能称为真实偏好的精确估计。

\begin{table}[htbp]
\centering
\caption{Global RRP 年度惩罚参数与验证结果}\label{tab:annual_penalties}
\resizebox{0.86\textwidth}{!}{\begin{tabular}{lrrrr}
\toprule
生效年 & $\lambda_{v,t}$ & $\lambda_{r,t}$ & 验证夏普 & 标准误集合\\\midrule
2017 & 6.858e-12 & 16.54 & 1.467 & 9/9\\
2018 & 0.1402 & 8.826 & 1.630 & 9/9\\
2019 & 0.08214 & 24.72 & 0.026 & 9/9\\
2020 & 0.1012 & 9.385 & 3.433 & 9/9\\
2021 & 0.1181 & 2.936 & 1.629 & 9/9\\
2022 & 0.08515 & 2.01 & 2.051 & 9/9\\
2023 & 0.1336 & 8.52 & 0.280 & 9/9\\
2024 & 0.2022 & 11.52 & 2.796 & 9/9\\
2025 & 0.1931 & 15.49 & 2.547 & 9/9\\
2026 & 0.1574 & 3.793 & 2.207 & 9/9\\
\bottomrule
\end{tabular}}
\end{table}

\subsection{凸性、可行性与唯一性}
跟踪项的 Hessian 为 $2I$，是严格凸函数。方差项在 $\Sigma_t\succeq0$ 且 $v_t>0$ 时为凸函数。收益短缺先对仿射函数取正部，再施加在非负域上单调的平方函数，因此保持凸性。所有尺度在求解前固定，正权重加总后目标仍然严格凸。

单纯形 $\Delta_t$ 是非空、闭且有界的凸集合，只要存在至少一个合格资产，等权向量便提供一个可行点。在代码进一步要求足够资产和历史观察的前提下，连续目标在该集合上存在最小值；严格凸性使最终权重解唯一。数值求解仍受精度限制，所以实际报告使用约束残差而非将机器输出当作精确实数解。

参考组合 $q_t$ 可以达到 $R_t$，所以收益基准在当期估计下可行。最终问题仍使用软惩罚，使权重能够在方差下降和预测收益短缺之间取舍。若最优权重的预测收益低于 $R_t$，差额会进入目标值，预算约束及非负约束不会因此放宽。

\subsection{零权重的最优性解释}
令 $\nu$ 为预算等式乘子，$\gamma_i\geq0$ 为非负约束乘子。最优性条件包括
\[
\nabla J_t(w)+\nu\mathbf1-\gamma=0,\qquad
\gamma_iw_i=0.
\]
对于正权重资产，$\gamma_i=0$，目标的边际变化在预算限制下平衡。对于零权重资产，非负约束可以生效，该资产在当期估计下增加资本未必改善目标。即使参考组合中的所有权重都为正，最终组合也可以合理地位于单纯形边界。

本文不对每只资产设置人为最低持仓。这样避免用极小权重制造“全部资产参与”的表面证据。资产参与由全部合格期间的真实目标权重记录判断，并同时查看持有期数和最大权重。永久不用与某期为零是不同问题，后文分别说明。

\subsection{数值求解与经济解释}
实际代码对矩阵进行对称化与特征值处理，使用支持凸锥规划的求解器。只有合规状态与有效变量同时满足要求时，才接受结果。失败不能静默替换为等权组合，否则发布的收益路径会混入没有标注的其他配置规则。

程序的 DCP 检查验证表达式结构，约束残差验证数值解。二者都不能证明预期收益准确，也不能将最大回撤锁定在历史观察的水平。凸性解决的是每一期给定输入下的分配问题，收益和风险表现还取决于输入估计及随后市场变化。本文将这一层次区别贯穿结果解释。
